% 2de graad, Driehoeksmeting (Pythagoras en de goniometrische getallen in rechthoekige driehoeken).
% Cézanne: de Mont Sainte-Victoire; een hoogte meet je met een hoek en een afstand.
% Motief: de rechthoekige driehoek 3-4-5, met de hoek alpha waarvoor tan alpha = 3/4.
\def\omslagbestand{omslag-sainte-victoire.jpg}
\def\omslagbron{Paul Cézanne, \textit{De Mont Sainte-Victoire}, ca. 1892 (publiek domein)}
\def\omslagrand{10.8 + 0.35*sin(\x*360/16)}
\def\omslagmotief{%
  \randlijn
  \begin{scope}[xshift=7mm]
  \draw[white,line width=1.4pt,line join=round] (14.6,4.8) -- (18.6,4.8) -- (18.6,7.8) -- cycle;
  \draw[white,line width=.8pt] (18.25,4.8) -- (18.25,5.15) -- (18.6,5.15);
  \draw[fel,line width=1.4pt] (15.3,4.8) arc[start angle=0,end angle=36.87,radius=.7];
  \node[text=fel,anchor=west] at (15.35,5.12) {$\alpha$};
  \node[text=white,font=\small,anchor=north] at (16.6,4.75) {$4$};
  \node[text=white,font=\small,anchor=west] at (18.65,6.3) {$3$};
  \node[text=white,font=\small,anchor=south east] at (16.55,6.35) {$5$};
  \node[text=white,anchor=south] at (16.6,8.2) {$3^2+4^2=5^2$\qquad$\tan\alpha=\frac34$};
  \end{scope}
}
